% GENERATED FILE. Edit canonical lessons, then run npm run generate:books.
% canonical-source-hash: fnv1a64:35cf749584aee856
% canonical-lessons: MR-C232-matsari, MR-C232-abhimani, MR-C232-asvastha, MR-C232-gulgulit, MR-C232-kharkharit

\chapter{Jealous, Proud, Uneasy, Smooth, Rough}
\label{ch:mr-jealous-proud-uneasy-smooth-rough}
\hlchaptermodality{232}

\begin{chapteropening}
\textbf{By the end of this chapter:} \emph{I can say jealous, proud, uneasy, smooth and rough.}

Complete the last lesson of chapter 232: I can say jealous, proud, uneasy, smooth and rough.
\end{chapteropening}

\section[\texorpdfstring{\mr{मत्सरी} (matsarī)}{मत्सरी (matsarī)}]{\mr{मत्सरी} (matsarī) — jealous}

\label{lesson:MR-C232-matsari}



\begin{quote}
Before the new one: say the Marathi for of course, then the Marathi for that is to say.
\end{quote}

\subsection*{You'll want to know: \mr{मत्सरी}}
\textbf{\mr{मत्सरी}} — \emph{matsarī} — \textquotedblleft{}jealous\textquotedblright{}.

\begin{cousinweb}
\textbf{\mr{मत्सरी}} (\emph{matsarī}) is \textbf{\mr{मत्सर}}, \textquotedblleft{}jealousy\textquotedblright{}, from Sanskrit \emph{matsara}, \textquotedblleft{}jealous\textquotedblright{} (CDIAL 9756), with the ending \emph{-ī}, \textquotedblleft{}having\textquotedblright{}.
\end{cousinweb}

\begin{grammarlens}[title={What you've built}]
One more word for this chapter.
\end{grammarlens}

\subsection*{Your turn}
\begin{itemize}
\raggedright
  \item \emph{Say it:} \emph{matsarī}
  \item \emph{Say it:} \emph{matsarī}, once more
  \item \emph{Say it:} say \emph{matsarī}
  \item \emph{Recall:} say the Marathi for of course, then the Marathi for that is to say, then say \emph{matsarī} again
\end{itemize}

\begin{tcolorbox}[breakable,colback=teal!4,colframe=teal!35!black,title={Before you move on}]
What does \mr{मत्सरी} mean? (\textquotedblleft{}Jealous\textquotedblright{}.) Say it once more.
\end{tcolorbox}
\section[\texorpdfstring{\mr{अभिमानी} (abhimānī)}{अभिमानी (abhimānī)}]{\mr{अभिमानी} (abhimānī) — proud}

\label{lesson:MR-C232-abhimani}



\begin{quote}
Before the new one: say the Marathi for that is to say, then the Marathi for jealous.
\end{quote}

\subsection*{You'll want to know: \mr{अभिमानी}}
\textbf{\mr{अभिमानी}} — \emph{abhimānī} — \textquotedblleft{}proud\textquotedblright{}.

\begin{cousinweb}
\textbf{\mr{अभिमानी}} (\emph{abhimānī}) is Sanskrit \emph{abhimānin}, \textquotedblleft{}proud\textquotedblright{}, a tatsama built on \emph{abhimāna}, \textquotedblleft{}pride\textquotedblright{} (CDIAL 524).
\end{cousinweb}

\begin{grammarlens}[title={What you've built}]
One more word for this chapter.
\end{grammarlens}

\subsection*{Your turn}
\begin{itemize}
\raggedright
  \item \emph{Say it:} \emph{abhimānī}
  \item \emph{Say it:} \emph{abhimānī}, once more
  \item \emph{Say it:} say \emph{abhimānī}
  \item \emph{Recall:} say the Marathi for that is to say, then the Marathi for jealous, then say \emph{abhimānī} again
\end{itemize}

\begin{tcolorbox}[breakable,colback=teal!4,colframe=teal!35!black,title={Before you move on}]
What does \mr{अभिमानी} mean? (\textquotedblleft{}Proud\textquotedblright{}.) Say it once more.
\end{tcolorbox}
\section[\texorpdfstring{\mr{अस्वस्थ} (asvastha)}{अस्वस्थ (asvastha)}]{\mr{अस्वस्थ} (asvastha) — uneasy, restless}

\label{lesson:MR-C232-asvastha}



\begin{quote}
Before the new one: say the Marathi for jealous, then the Marathi for proud.
\end{quote}

\subsection*{You'll want to know: \mr{अस्वस्थ}}
\textbf{\mr{अस्वस्थ}} — \emph{asvastha} — \textquotedblleft{}uneasy, restless\textquotedblright{}.

\begin{cousinweb}
\textbf{\mr{अस्वस्थ}} (\emph{asvastha}) is Sanskrit \emph{asvastha}, \textquotedblleft{}not at ease\textquotedblright{}, a tatsama: \emph{a-}, \textquotedblleft{}not\textquotedblright{}, and \emph{svastha}, \textquotedblleft{}well, healthy\textquotedblright{} (CDIAL 13917).
\end{cousinweb}

\begin{grammarlens}[title={What you've built}]
One more word for this chapter.
\end{grammarlens}

\subsection*{Your turn}
\begin{itemize}
\raggedright
  \item \emph{Say it:} \emph{asvastha}
  \item \emph{Say it:} \emph{asvastha}, once more
  \item \emph{Say it:} say \emph{asvastha}
  \item \emph{Recall:} say the Marathi for jealous, then the Marathi for proud, then say \emph{asvastha} again
\end{itemize}

\begin{tcolorbox}[breakable,colback=teal!4,colframe=teal!35!black,title={Before you move on}]
What does \mr{अस्वस्थ} mean? (\textquotedblleft{}Uneasy, restless\textquotedblright{}.) Say it once more.
\end{tcolorbox}
\section[\texorpdfstring{\mr{गुळगुळीत} (guḷguḷīt)}{गुळगुळीत (guḷguḷīt)}]{\mr{गुळगुळीत} (guḷguḷīt) — smooth}

\label{lesson:MR-C232-gulgulit}



\begin{quote}
Before the new one: say the Marathi for proud, then the Marathi for uneasy.
\end{quote}

\subsection*{You'll want to know: \mr{गुळगुळीत}}
\textbf{\mr{गुळगुळीत}} — \emph{guḷguḷīt} — \textquotedblleft{}smooth\textquotedblright{}.

\begin{cousinweb}
The origin of \textbf{\mr{गुळगुळीत}} (\emph{guḷguḷīt}) is not settled. Neither CDIAL nor DEDR, searched by form and by meaning, has an entry for it.
\end{cousinweb}

\begin{grammarlens}[title={What you've built}]
One more word for this chapter.
\end{grammarlens}

\subsection*{Your turn}
\begin{itemize}
\raggedright
  \item \emph{Say it:} \emph{guḷguḷīt}
  \item \emph{Say it:} \emph{guḷguḷīt}, once more
  \item \emph{Say it:} say \emph{guḷguḷīt}
  \item \emph{Recall:} say the Marathi for proud, then the Marathi for uneasy, then say \emph{guḷguḷīt} again
\end{itemize}

\begin{tcolorbox}[breakable,colback=teal!4,colframe=teal!35!black,title={Before you move on}]
What does \mr{गुळगुळीत} mean? (\textquotedblleft{}Smooth\textquotedblright{}.) Say it once more.
\end{tcolorbox}
\section[\texorpdfstring{\mr{खरखरीत} (kharkharīt)}{खरखरीत (kharkharīt)}]{\mr{खरखरीत} (kharkharīt) — rough}

\label{lesson:MR-C232-kharkharit}



\begin{quote}
Before the new one: say the Marathi for uneasy, then the Marathi for smooth.
\end{quote}

\subsection*{You'll want to know: \mr{खरखरीत}}
\textbf{\mr{खरखरीत}} — \emph{kharkharīt} — \textquotedblleft{}rough\textquotedblright{}.

\begin{cousinweb}
The origin of \textbf{\mr{खरखरीत}} (\emph{kharkharīt}) is not settled. Neither CDIAL nor DEDR, searched by form and by meaning, has an entry for it.
\end{cousinweb}

\begin{grammarlens}[title={What you've built}]
One more word for this chapter.
\end{grammarlens}

\subsection*{Your turn}
\begin{itemize}
\raggedright
  \item \emph{Say it:} \emph{kharkharīt}
  \item \emph{Say it:} \emph{kharkharīt}, once more
  \item \emph{Say it:} say \emph{kharkharīt}
  \item \emph{Recall:} say the Marathi for uneasy, then the Marathi for smooth, then say \emph{kharkharīt} again
\end{itemize}

\begin{tcolorbox}[breakable,colback=teal!4,colframe=teal!35!black,title={Before you move on}]
What does \mr{खरखरीत} mean? (\textquotedblleft{}Rough\textquotedblright{}.) Say it once more.
\end{tcolorbox}
